\section{Computing the moments and designing the experiments}\label{sec:methods}
\subsection{A deterministic evaluator}
The statistical moment formulas are exact; their implementation evaluates
the required input integrals numerically. For a density $f$ and rounding
cell $J_{d,k}=[(d-1/2)s_k,(d+1/2)s_k)$, the continuous contribution is
\begin{equation}\label{eq:cell-integration}
 T_{\ell,k}(t)=\sum_d\int_{J_{d,k}\cap(-t,t)}
                         (d s_k-x)^\ell f(x)\,dx.
\end{equation}
Atoms contribute their probabilities times $h_k(x)^\ell$, with ties to
even. Expanding the integrand uses cell probabilities and the first two
truncated input moments. Since $t/s_k\le2^{L+1}$ at both binade endpoints,
each exponent requires $O(2^L)$ cells, independently of $n$.
For heterogeneous inputs, these integrals are evaluated coordinatewise.
The exponential identity \eqref{eq:revision-exp-closed} supplies a separate
closed-form check.

We form $G(t)^n$ from survival probabilities and logarithms, and use
\texttt{expm1} differences for nearby log probabilities.
Equation~\eqref{eq:revision-negative} evaluates the conditional covariance
stably; centered sums across exponent states evaluate $B$. The evaluator
records omitted exponent probability and numerical stability diagnostics.
The large-block exponential calculation uses higher precision to resolve
survival probabilities below binary64 resolution.

We compare four predictions. The \emph{finite-block prediction} uses
\eqref{eq:revision-global}--\eqref{eq:revision-decomposition}.
The \emph{conditional-mean upper prediction} retains $m_k,t_k,v$ and sets
$C=0$, as in Corollary~\ref{cor:revision-bracket}.
The \emph{unconditioned-response approximation} mixes moments of $h_k(X)$
under the original input law with $P_k$ and assumes conditional
independence. The \emph{uniform-noise prediction} uses
$\mu=\kappa=0$ and $v=\mathbb E[s_{K_n}^2]/12$.

For an iid law, the evaluator visits the retained exponent states and
integrates $O(2^L)$ rounding cells per state, so its scalar work is
$O(K2^L)$ for $K$ retained states and does not grow with the block length
$n$. For a heterogeneous block of size $b$, the corresponding work is
$O(K(b2^L+b^2))$ including covariance assembly; storage is $O(b^2)$.
This is the cost of the population evaluator, separate from Monte Carlo
sampling and from the linear solve used in the application experiments.

In the tables, ``Exact'' denotes numerical evaluation of the exact moment
identities. The heterogeneous and latent evaluators retain exponents
$-40\le k\le35$ and skip states of probability below $10^{-25}$.
They record omitted-probability and residual-energy bounds separately from
integration convergence checks. Numerical truncation bounds are therefore
part of the population prediction, rather than evidence supplied by its
Monte Carlo interval.
The application codecs use an 8-bit signed exponent and reject out-of-range
states; all retained source vectors and all population states lie inside that
range, so no saturation or exponent clipping is included in the reported
errors.

\subsection{Independent verification}
An independent rational oracle enumerates $25{,}600$ blocks in $96$
configurations: six finite-support laws,
$n\in\{2,3,4,6\}$, and $L\in\{0,1,2,4\}$.
It checks the conditional moments, covariance eigenvalues, weighted MSE,
and certified bracket, including atoms, zero blocks, dyadic endpoints,
and ties to even. A second oracle enumerates $480$ blocks in $32$
heterogeneous-input configurations and checks
Proposition~\ref{prop:revision-nonidentical}, including both conditional
covariance signs. Rational Green-matrix calculations verify linear
propagation against direct gradient energy in one and two dimensions.
Independent lognormal evaluators use analytical cell moments at $60$ and
$90$ decimal digits to check the homogeneous and heterogeneous moment
calculations. An $80$-digit Pareto cell integrator agrees with all eight
binary64 inflation factors to relative error below $8.1\times10^{-10}$.

\subsection{Sampling and measurement}
The main-grid and weighted experiments use independent blocks, master seed
$20261010$, and labeled child random streams. Population predictions are
computed from the input law. For each sampled block we retain
$S_b=\sum_i e_{b,i}$, $Q_b=\sum_i e_{b,i}^2$, exponent state, and squared
step, with $s_*=0$. For $B_{\rm MC}$ blocks,
\begin{equation}\label{eq:mc-estimator}
 \widehat\mu=\frac{\sum_b S_b}{nB_{\rm MC}},\quad
 \widehat v=\frac{\sum_b Q_b}{nB_{\rm MC}}-\widehat\mu^2,\quad
 \widehat{\mathcal I}_n=
   \frac{\operatorname{sampleVar}_b(S_b)}{n\widehat v}.
\end{equation}
The numerator has denominator $B_{\rm MC}-1$. Confidence intervals are
nominal $95\%$ percentile intervals from $200$ bootstrap resamples of whole
blocks. We use ``contains'' below to report whether a displayed interval
contains the population prediction; this is an empirical interval check,
not a repeated-sampling coverage theorem. Application-specific streams,
sample counts, and vector-bootstrap protocols are given with their
experiments.

The main grid uses $\mathcal N(0,1)$, $\exp(Z)$, $|Z|$, $\max(Z,0)$
for $Z\sim\mathcal N(0,1)$, $\operatorname{Exp}(1)$, and Pareto$(4)$
with $\mathbb P(X>x)=x^{-4}$ for $x\ge1$.
Combining $n\in\{32,128,512,4096\}$ with $L\in\{2,4\}$ gives $48$
round-to-nearest configurations, with respectively
$100{,}000$, $40{,}000$, $20{,}000$, and $6{,}000$ blocks.
Eight stochastic-rounding controls use $n=128$, both precisions, and
Gaussian, lognormal, half-normal, and ReLU inputs.
Six small-block cases use Beta$(2,5)$ and exponential inputs at $n=2,4,8$,
$L=2$, with $150{,}000$ blocks each; two more Beta cases use
$n=32,128$, $L=2$, and the corresponding main-grid block counts.
This gives $64$ configurations.

Pareto$(4)$ has an infinite fourth moment, and ordinary samples can miss
rare exponent states that contribute substantially to second moments.
We therefore sample all eight Pareto configurations by exponent strata
and retain the ordinary results in the package. State $k$ receives
$\max(100,\operatorname{round}(30{,}000P_k))$ conditional blocks.
To draw one, sample the maximum by inversion uniformly between
$G(a_k)^n$ and $G(b_k)^n$, then draw the remaining $n-1$ coordinates
independently below that maximum. Placing the maximum last preserves
the symmetric statistics $S_b,Q_b$.

The stratified estimator combines the first and second moments of
$S_b,Q_b$ with analytical state probabilities, renormalized over retained
states. A variance correction removes the squared-mean estimation bias.
Bootstrap resampling occurs within each stratum with its population
weight fixed. Allocations and omitted-mass diagnostics are saved.
For Pareto$(4)$ and $T\ge1$, the omitted upper-tail marginal second moment
satisfies
\begin{equation}\label{eq:pareto-tail-bound}
 \mathbb E[e_1^2\mathbf1_{\{M_n\ge T\}}]
 \le\frac{\mathbb E[M_n^2\mathbf1_{\{M_n\ge T\}}]}{4\,2^{2L}}
 \le\frac{n}{2\,2^{2L}T^2}.
\end{equation}

Weighted reductions use $50{,}000$ paired lognormal blocks at $n=128$,
$L=2$, deterministic coefficients, and binary64 products and reductions.
We record mean error, centered variance, and MSE separately for shared
scales and partitions. The exponential asymptotic experiment evaluates
the moment formulas at high precision along the theorem's subsequence.
